\begin{textsample}{POS Dim 6 – ma – Score 8.00 – ma/ma\_em\_44.txt}  \label{ex:f6_pos_097}
Proposição de itinerários \textbf{formativos} integrados Marco legal Conceitualmente, os itinerários \textbf{formativos} são o conjunto de unidades curriculares ofertadas pelas instituições das redes de ensino que possibilitam ao estudante aprofundar seus conhecimentos e se preparar para o prosseguimento de estudos ou para o mundo do trabalho, de forma a contribuir para a construção de soluções de \textbf{problemas} específicos da sociedade.
As unidades curriculares, por sua vez, são elementos com carga horária pré-definida, formadas pelo conjunto de estratégias cujo objetivo é desenvolver competências específicas ( Resolução CNE nº 3/2018 ).
As bases legais para organização curricular e metodológica dos itinerários \textbf{formativos} estão previstos na Lei nº 13.415/2017, na Resolução do CNE nº 3/2018 e na Portaria MEC nº 1.432/2018, que estabelece os referenciais para elaboração dos itinerários \textbf{formativos}.
Os itinerários \textbf{formativos} compõem a parte diversificada do currículo, que também é formada pelos componentes projeto de vida, eletivas de base, eletivas de pré-itinerário \textbf{formativo}, tutoria e corresponsabilidade social, elementos obrigatórios que, juntamente com a formação geral básica, fazem parte da trilha/arranjo curricular nas três séries do ensino médio para os estudantes.
A rede estadual de ensino do estado do Maranhão oferece as 1.800 horas de BNCC distribuídas nos três anos, garantidas as previsões legais, como o oferecimento de língua portuguesa e matemática nos três anos, e língua inglesa como obrigatória ( Lei nº 13.415/17 ).
Assim, para os itinerários \textbf{formativos}, a rede dispõe de 1.200 horas distribuídas nas três séries do ensino médio, observadas as particularidades da escolha de cada estudante, em acordo ao seu projeto de vida, e estão assim distribuídos : - Lei nº 13.415 ( 16/2/2017 ) - DCNEM ( 8/11/2018 ) - BNCC Ensino Médio ( 4/12/2018 ) - Referenciais dos itinerários \textbf{formativos} ( 28/12/2018 ) Em acordo à Resolução nº 3, de 21 de novembro de 2018 ( DCNEM ), > Art. 12.
A partir das áreas do conhecimento e da formação técnica e profissional, os itinerários \textbf{formativos} devem ser organizados, considerando : I - linguagens e suas \textbf{tecnologias} : aprofundamento de conhecimentos \textbf{estruturantes} para aplicação de diferentes linguagens em contextos sociais e de trabalho, estruturando arranjos curriculares que permitam estudos em línguas vernáculas, estrangeiras, clássicas e indígenas, Língua Brasileira de Sinais ( LIBRAS ), das artes, design, linguagens \textbf{digitais}, corporeidade, artes cênicas, roteiros, produções literárias, entre outros, considerando o contexto local e as possibilidades de oferta pelos sistemas de ensino ; > > IV - ciências humanas e sociais aplicadas : aprofundamento de conhecimentos \textbf{estruturantes} para aplicação de diferentes conceitos em contextos sociais e de trabalho, estruturando arranjos curriculares que permitam estudos em relações sociais, modelos econômicos, processos políticos, pluralidade cultural, historicidade do universo, do homem e natureza, entre outros, considerando o contexto local e as possibilidades de oferta pelos sistemas de ensino.
Objetivos Os objetivos dos itinerários \textbf{formativos} percorrem a trajetória de : - Aprofundar as aprendizagens relacionadas às competências gerais, às áreas de conhecimento e/ou à formação técnica e profissional ; - Consolidar a formação integral dos estudantes, desenvolvendo a autonomia necessária para que realizem seus projetos de vida ; - Promover a incorporação de valores universais, como ética, liberdade, democracia, justiça social, pluralidade, solidariedade e \textbf{sustentabilidade} ; - Desenvolver habilidades que permitam aos estudantes ter uma visão de mundo ampla e heterogênea, tomar decisões e agir nas mais diversas situações, seja na escola, seja no trabalho, seja na vida.
Os itinerários \textbf{formativos} trazem a possibilidade da diversificação curricular, já que o estudante terá autonomia para escolha de seu percurso \textbf{formativo}.
Para isso, distintas trajetórias serão oferecidas ao estudante para serem desenvolvidas por meio de diferentes desenhos curriculares.
Isso tudo contribuindo para aprofundar e/ou ampliar as aprendizagens do estudante em uma ou mais áreas de conhecimento, por meio da aquisição das habilidades referenciadas nos eixos \textbf{estruturantes}, tendo por objetivo criar as condições necessárias para que ele seja formado para prosseguir nos seus estudos e/ou inserção no mundo do trabalho, executando o seu projeto de vida.
Os itinerários \textbf{formativos} como forma de aprofundamentos dos conhecimentos, além de melhor explorar potenciais e vocações dos estudantes, permitem que os jovens já concluam o ensino médio com algum diferencial na sua formação.
Desta forma, o histórico escolar será personalizado para atender à peculiaridade e maior tempo de dedicação às unidades curriculares escolhidas pelo estudante conjugadas com seu projeto de vida.
Para que esses objetivos sejam cumpridos, a orientação é que os aprofundamentos/itinerários \textbf{formativos} tenham duração de, pelo menos, quatro semestres ( Consed a, 2020 ).
Nessa perspectiva, a escolha do itinerário \textbf{formativo} se dará ao final da 1ª série, o estudante cursará seu itinerário a partir da 2ª série e estará associado a outras unidades curriculares, conforme o itinerário escolhido.
Ressalta -se que os itinerários \textbf{formativos} se constituem em campos de conhecimentos que carregam uma abordagem inter e transdisciplinar e devem se estruturar a partir de percurso com começo, meio e fim, cujo fluxo atravesse os quatro eixos \textbf{estruturantes} ( investigação científica, \textbf{mediação} e intervenção sociocultural, processos criativos e \textbf{empreendedorismo} ) e permita aos estudantes se desenvolverem de forma integral, orgânica e progressiva, lidando adequadamente com desafios cada vez mais complexos.
Também é interessante que cada etapa dessa jornada integre e articule os conhecimentos, habilidades, atitudes e valores adquiridos nas etapas anteriores.
O compromisso da educação por meio do desenvolvimento humano global, em suas dimensões intelectual, física, afetiva, social, ética, moral e simbólica, são prerrogativas já consolidadas nos marcos legais como a LDB, a BNCC e as DCNs para o ensino médio.
Outro aspecto fundamental para os currículos refere -se à prática pedagógica orientada sob a perspectiva de organização interdisciplinar e transdisciplinar dos componentes curriculares.
Tal prática deve ser planejada sob a perspectiva do aprofundamento de objetos de conhecimento, por meio de seus conceitos e compreensão das correlações que envolvem os conteúdos específicos de cada componente curricular.
Ainda segundo o artigo 7º, parágrafo 2º das DCNs : > O currículo deve contemplar tratamento metodológico que evidencie a contextualização, a diversificação e a transdisciplinaridade ou outras formas de interação e articulação entre diferentes campos de saberes específicos, contemplando vivências práticas e vinculando a educação escolar ao mundo do trabalho e à prática social e possibilitando o aproveitamento de estudos e o reconhecimento de saberes adquiridos nas experiências pessoais, sociais e do trabalho ( DCN, art. 7, § 2º ).
Desse modo, ao articular os itinerários \textbf{formativos} de forma transdisciplinar, entende -se que a organização e o direcionamento da prática estejam voltados para realização de diálogos entre os componentes, tornando o currículo cada vez mais próximo do estudante protagonista, ao ser tratado de forma contextualizada e interdisciplinar ( art. 11, § 5º ).
Essa perspectiva de ampliação de possibilidades que o itinerário \textbf{formativo} pode oferecer corrobora o reflexo que encontramos atualmente na escola – uma diversificação e, partindo desse entendimento, compreendemos que as aprendizagens dos estudantes precisam se desenvolver nesse caminho.
Assim, torna -se interessante considerarmos o pensamento de Zabala ( 2002, p. 15 ) : “ O reflexo, na escola, dessa diversificação é a seleção ou a distribuição dos conteúdos escolares a partir de parâmetros basicamente disciplinares ”.
Nessa mesma direção, Edgar Morin ( 2013 ) trabalha a ideia de um mundo globalizado e acelerado.
Nesse sentido, compreendemos que esse pensamento não pode passar despercebido para o campo da educação.
Assim, defende a importância de transmitir ao estudante que ele é um ser multidimensional.
Morin ( 2017 ) defende, também, a unidade do conhecimento, pois acredita que as disciplinas fechadas impedem a compreensão dos \textbf{problemas} do mundo.
Nessa perspectiva, a transdisciplinaridade é fundamental para englobar e garantir a compreensão do saber.
Para o autor, é justamente a transdisciplinaridade que possibilita, por meio das disciplinas, a transmissão de uma visão de mundo mais complexa.
Entendemos, portanto, que o currículo do estado precisa caminhar junto com o mundo globalizado e que atenda de fato às reais necessidades dos estudantes.
Dessa forma, acreditamos ser fundamental considerarmos em nosso currículo essa visão interdisciplinar¹⁴, transdisciplinar¹⁵ e globalizadora na composição dos itinerários \textbf{formativos}.
Nessa perspectiva, a rede estadual de ensino do Maranhão adotará a proposta de itinerários \textbf{formativos} integrados por meio de aprofundamentos que perpassam todas as áreas de conhecimentos, com enfoque globalizador, em uma perspectiva interdisciplinar e transdisciplinar, visto que os estudantes do ensino médio, em geral, concluem essa etapa com a visão na perspectiva do Enem.
Vale destacar que os itinerários \textbf{formativos} integrados, nessa proposta, apresentam -se em torno de arranjos curriculares que vão atender às características peculiares e específicas dos cursos das etapas subsequentes ao ensino médio, por meio da flexibilização e integração das áreas de conhecimento, agrupados por afinidades, similaridades e atributos comuns dos cursos ofertados nas etapas de ensino superior e serão desenvolvidos por meio de eletivas definidas e elaboradas por cada estabelecimento de ensino.
É importante ressaltar que esses arranjos curriculares para a proposta dos itinerários \textbf{formativos} integrados foram reorganizados a partir de pesquisas e levantamentos realizados pela equipe do Currículo da Seduc -MA.
Essas observações foram realizadas nas instituições de ensino : UFMA, UEMA, IFMA e Ceuma.
Como resultado, observou -se que as universidades do país e, sobretudo, as do estado do Maranhão têm autonomia para organizar os cursos em centros, de acordo com as áreas afins.
Nessa perspectiva, os cursos de graduação possuem seus desdobramentos para a mobilidade e empregabilidade dos egressos de cada curso.
Assim, as nomenclaturas são definidas de acordo com suas realidades e demandas locais.
Foi a partir da organização das instituições do estado do Maranhão que os arranjos curriculares foram agrupados e pensados na proposta de cada itinerário integrado que será apresentada.
Cabe destacar, ainda, que, segundo o que é estabelecido pelo MEC, os centros universitários são instituições pluricurriculares, abrangendo uma ou mais áreas de conhecimento.
Para organização dos itinerários \textbf{formativos} integrados do estado do Maranhão, é necessário considerar a nossa realidade, uma vez que o arranjo para os itinerários vai priorizar especificamente os estudantes locais.
É preciso, sobretudo, olhar para a diversidade e pluralidade desse universo que é a escola.
Assim, as aplicabilidades desses itinerários dependem das possibilidades de articulação de cada município no que se refere aos custos e à própria estrutura da escola e da localidade.
Com base nisso, a partir do levantamento realizado nos cursos ofertados nas instituições de ensino citadas, tivemos a composição dos arranjos curriculares distribuídos em quatro itinerários \textbf{formativos} integrados, que serão ofertados na rede estadual de ensino, a saber : i ) Ciências Exatas, Tecnológicas e da Terra ; ii ) Ciências da Saúde ; iii ) Ciências Humanas e Linguagens ; iv ) Ciências Sociais, Econômicas e Administrativas.
Considerando essa organização dos itinerários \textbf{formativos} integrados, deve -se levar em consideração que, para a construção das eletivas dos itinerários \textbf{formativos}, além de considerar o contexto local, é preciso assegurar padrões básicos de aprendizagem e de ensino que estarão garantidos na articulação com os aprofundamentos dos conhecimentos dos componentes curriculares que compõem cada itinerário x competências x habilidades.
Tais aspectos também se articulam com a proposta pedagógica da escola, que estarão refletidos nos planos das atividades docentes. ¹⁴ A interdisciplinaridade propõe a capacidade de dialogar entre as diversas ciências, fazendo entender o saber como um todo, e não como partes fragmentadas ( MORIN, 2005 ). ¹⁵ A ideia de transdisciplinaridade surgiu para superar o conceito de disciplina, que se configura pela departamentalização do saber em diversas matérias, em que cada disciplina é abordada de modo fragmentado e isolada das demais.
A transdisciplinaridade não significa apenas que as disciplinas colaboram entre si, mas, também, que existe um pensamento organizador que ultrapassa as próprias disciplinas.
Para haver essa transdisciplinaridade, é preciso haver um pensamento organizador, chamado pensamento complexo.

% matched lemmas: digital, empreendedorismo, estruturante, formativo, mediação, problema, sustentabilidade, tecnologia
\end{textsample}
